\documentclass[11pt,reqno]{amsart}

%=====================================================================
%  Team A1 -- Role 1 (Optimal Contract and General Carbon Theory)
%  A1 WEEK 4.  The carbon classification theorem.
%   1. A single closed-form functional Phi(F) with sign(Delta)=sign(Phi) -- both
%      directions (an exact iff), from the Week-2 affine representation.
%   2. Minimality of the statistic: both conditional means are needed (the sharp
%      converse to Week-2's not-cutoff-decidable impossibility).
%   3. The frontier leverage condition, closed toward a theorem (exact reduction
%      + a proved K_F-monotonicity lemma + the verified scalar).
%   4. The carbon phase diagram: sign(Phi) partitions the primitive space.
%=====================================================================

\usepackage[T1]{fontenc}
\usepackage{lmodern}
\usepackage{amsmath,amssymb,amsthm,mathtools}
\usepackage[margin=1in]{geometry}
\usepackage{microtype}
\usepackage{enumitem}
\usepackage{booktabs}
\usepackage{xcolor}
\usepackage[colorlinks=true,linkcolor=blue!55!black,citecolor=blue!55!black]{hyperref}

\theoremstyle{plain}
\newtheorem{theorem}{Theorem}
\newtheorem{proposition}{Proposition}
\newtheorem{lemma}{Lemma}
\newtheorem{corollary}{Corollary}
\theoremstyle{definition}
\newtheorem{condition}{Condition}
\newtheorem{assumption}{Assumption}
\newtheorem*{cardef}{Definition}
\theoremstyle{remark}
\newtheorem{remark}{Remark}

\newcommand{\KR}{K_R}\newcommand{\KF}{K_F}\newcommand{\KM}{K_M}
\newcommand{\KRs}{K_R^{*}}\newcommand{\KFs}{K_F^{*}}\newcommand{\KMs}{K_M^{*}}
\newcommand{\KRo}{K_R^{\circ}}\newcommand{\KFo}{K_F^{\circ}}
\newcommand{\Dbar}{\bar D}\newcommand{\CM}{C_M}\newcommand{\CF}{C_F}\newcommand{\CR}{C_R}
\newcommand{\bp}{\beta+p}\newcommand{\FKR}{F_{\!\KR}}
\newcommand{\E}{\mathbb{E}}\newcommand{\PP}{\mathbb{P}}\newcommand{\eqdef}{\vcentcolon=}
\newcommand{\Ucal}{\mathcal U}\newcommand{\fossil}{\mathrm{fossil}}
\newcommand{\mustar}{\mu^{*}_{\mathrm{low}}}\newcommand{\mumid}{\mu^{\circ}_{\mathrm{mid}}}
\newcommand{\phio}{\varphi^{\circ}}\newcommand{\qo}{q^{\circ}}

\begin{document}

\title[The Carbon Classification Theorem (Week 4)]
{The Carbon Classification Theorem:\\
An Exact Closed-Form Criterion for When a Flexible Load Reduces Emissions}
\author{Team A1 --- Role 1 (Optimal Contract and General Carbon Theory)}
\date{Working results, Week~4}

\begin{abstract}
This is the Role~1 Week~4 deliverable: an exact, computable criterion for the carbon
sign. We construct a carbon functional $\Phi(F)$---an explicit affine combination of the
distribution-free cutoffs, the equilibrium capacities, the surplus probability
$\qo=F(0;\KRo)$, and exactly two conditional means of the residual-demand law---and show
$\Delta=\Phi(F)$, so carbon falls if and only if
$\Phi(F)>0$. The identity is Week-2 bookkeeping; the value is that $\Phi$ is a
\emph{closed form on a named, minimal set of coordinates}, and that the sign needs
distributional information no coarser than those coordinates: both conditional means enter
with nonzero weight, and (Week~2) the sign is not decidable from capacities and cutoffs
alone over the admissible-law class. We are explicit about scope. $\Phi$ is the criterion
for the \emph{naive} dispatch of Weeks~1--2, for which---by the team's own results---%
$\Phi<0$ at interior optima (carbon rises) and $\Phi>0$ only at the cheap-renewable
corner. The point of naming $\Phi$ is exactly to \emph{explain the Week-2$\to$Week-3
flip}: the surplus contract of Week~3 zeroes the utilization term and turns $\Phi$
positive, so the classification is the ledger on which ``why the optimal contract
decarbonizes'' is read off. The three questions that a scan alone had settled are now \emph{conditional theorems}, each
gated on one explicit, checkable condition that we then verify: (a) each mean is
individually sign-decisive \emph{iff} a scalar margin beats the carbon deficit
($\mathrm{D}_\ell,\mathrm{D}_m$)---an exact both-directions criterion; (b) the carbon-down
set is an interval under monotonicity of $\Phi(\CR)$ (Condition~S); and (c) surplus-only is
globally carbon-optimal ($\Lambda\le1$) under a transparent capacity-monotonicity condition
(the budget raises fossil, not renewable, capacity---Condition~M), which we \emph{prove}
forces $\Lambda\le0$. Each condition holds throughout the scan. This converts three
numerical claims into proved statements whose only ungoverned inputs are named, checkable
inequalities on equilibrium quantities. Every identity is checked against quadrature in
\texttt{verify\_week4.py}: $\Phi$ matches $\Delta$ to machine precision, the margins
$\mathrm{D}_\ell,\mathrm{D}_m$ and Conditions~S,~M all hold, and the per-mean \emph{iff} is
confirmed exact.
\end{abstract}

\maketitle

%---------------------------------------------------------------------
\section{Setup: the affine representation}\label{sec:setup}

We keep the model and notation of Weeks~1--3. A regulated utility chooses renewable
$\KR$, fossil $\KF$, and flexible-computing-load (\textsc{fcl}) $\KM$ capacity; residual
demand is $\tilde x=\tilde D-\KR\tilde S\in[-\KR,\Dbar]$ with CDF $F(x;\KR)$, density
$f>0$ on the interior, $\FKR\eqdef\partial F/\partial\KR>0$. Assumption~1
($R>c+\CM$, $R<c+\CF+\CM$, $\bp>c+\CF$, $\bp>R$) holds throughout. The
distribution-free cutoffs are
$\varphi_1=F(a;\KRs)=\CM/(R-c)$, $\varphi_2=F(\KFs;\KRs)=(\bp-c-\CF-\CM)/(\bp-R)$, and
$\phio=F(\KFo;\KRo)=(\bp-c-\CF)/(\bp-c)$, with $a\eqdef\KFs-\KMs$.

The keystone of Week~2 is that net fossil is \emph{affine} in two conditional means.
Recall (Week~2, eq.~(fos)) the exact identities
\begin{equation}\label{eq:fos}
  \E[\fossil^{*}]=\KFs-\tfrac{\CM}{R-c}\big(a-\mustar\big),\qquad
  \E[\fossil^{\circ}]=\KFo(1-\phio)+(\phio-\qo)\,\mumid,
\end{equation}
where $\qo=F(0;\KRo)$ and
\begin{equation}\label{eq:means}
  \mustar\eqdef\E[\tilde x\mid\tilde x\in(-\KRs,a)]\in(-\KRs,a),\qquad
  \mumid\eqdef\E[\tilde x\mid\tilde x\in(0,\KFo)]\in(0,\KFo)
\end{equation}
are the conditional means over the \textsc{fcl} surplus zone and the no-\textsc{fcl}
mid zone. The net carbon saving is $\Delta\eqdef\E[\fossil^{\circ}]-\E[\fossil^{*}]$;
carbon falls iff $\Delta>0$.

%=====================================================================
\section{The carbon classification theorem}\label{sec:classify}
%=====================================================================

\begin{cardef}[The carbon functional]
At the equilibrium define
\begin{equation}\label{eq:Phi}
  \Phi(F)\;\eqdef\;
  \underbrace{\tfrac{\CM}{R-c}\big(a-\mustar\big)}_{\text{capacity substitution}}
  \;-\;
  \underbrace{\Big[\KFs-\KFo(1-\phio)-(\phio-\qo)\,\mumid\Big]}_{\text{net utilization}} .
\end{equation}
$\Phi$ is a closed-form affine functional of the statistic
$T(F)\eqdef\big(\KFs,\KMs,\KRs,\KFo,\phio,\qo,\mustar,\mumid\big)$: distribution-free
cutoffs, the equilibrium capacities, and the two conditional means \eqref{eq:means}.
\end{cardef}

\begin{theorem}[Carbon classification]\label{thm:classify}
For every residual-demand law $F$ with an interior optimum,
\[
  \boxed{\;\Delta=\Phi(F),\qquad\text{hence}\qquad \text{carbon falls}\iff \Phi(F)>0.\;}
\]
The equivalence is exact and holds in both directions; $\Phi>0$ is necessary and
sufficient for $\Delta>0$, $\Phi=0$ for carbon-neutrality, and $\Phi<0$ for carbon
increase.
\end{theorem}

\begin{proof}
Subtract the two identities in \eqref{eq:fos}:
\[
  \Delta=\E[\fossil^{\circ}]-\E[\fossil^{*}]
  =\KFo(1-\phio)+(\phio-\qo)\mumid-\KFs+\tfrac{\CM}{R-c}(a-\mustar),
\]
which is exactly $\Phi(F)$ of \eqref{eq:Phi}. As $\Delta=\Phi$ identically, every sign
statement transfers. Both directions are therefore immediate: $\Delta>0\Rightarrow\Phi>0$
and $\Phi>0\Rightarrow\Delta>0$.\footnote{Verified in \texttt{verify\_week4.py}: the
closed-form $\Phi$ of \eqref{eq:Phi} equals the direct quadrature of
$\E[\fossil^{\circ}]-\E[\fossil^{*}]$ to $\sim10^{-9}$ at every optimum across the
distribution library, and $\operatorname{sign}\Delta=\operatorname{sign}\Phi$ in every
cell.}
\end{proof}

The content of Theorem~\ref{thm:classify} is not the identity---which is Week-2
bookkeeping---but that the carbon sign is a \emph{closed form on a named, small set of
coordinates}: the distribution-free cutoff $\phio$, four capacities
$\KFs,\KMs,\KRs,\KFo$ (equivalently the cutoffs $\varphi_1,\varphi_2$ and $\phio$), the
surplus probability $\qo=F(0;\KRo)$, and the two conditional means. Of these only $\phio$
(and $\varphi_1,\varphi_2$) is distribution-free; the capacities, $\qo$, and the means are
all functionals of $F$---so $\Phi$ depends on $F$ through both the equilibrium map and the
two moments, and ``coordinates'' below means these entries, not a probabilistic sufficient
statistic. Whether the sign reduces further---to the distribution-free cutoffs alone---is
the subject of the next section, and the answer is no.

\begin{corollary}[The sufficient conditions, recovered]\label{cor:suff}
The Week-1/2 sufficient condition follows by bounding the means:
$\mustar\le(a-\KRs)/2$ and $\mumid\ge\KFo/2$ (monotone-density shape class $\Ucal$) give
$\Phi\ge\Phi_0\eqdef\tfrac{\CM}{R-c}\tfrac{a+\KRs}{2}-\KFs+\KFo[(1-\phio)+\tfrac{\phio-\qo}{2}]$,
so $\Phi_0>0$ implies carbon falls. The classification subsumes it: the sufficient
condition is the midpoint slice of the exact criterion $\Phi>0$.
\end{corollary}

%=====================================================================
\section{Minimality: the sign needs both means}\label{sec:minimal}
%=====================================================================

Week~2 proved the carbon sign is not a function of the distribution-free cutoffs
alone: over the class of admissible conditional laws sharing the equilibrium cutoffs and
capacities, $\Delta$ ranges over an interval straddling zero. We record what this says
about $\Phi$'s arguments.

\begin{proposition}[Both means are load-bearing]\label{thm:minimal}
As a function of $(\mustar,\mumid)$,
\[
  \frac{\partial\Phi}{\partial\mustar}=-\frac{\CM}{R-c}<0,\qquad
  \frac{\partial\Phi}{\partial\mumid}=\phio-\qo>0 ,
\]
both nonzero: $\Phi$ depends nontrivially on each conditional mean, so neither can be
dropped from the criterion. Consequently the distribution-free cutoffs and the capacities
do \emph{not} determine $\operatorname{sign}\Delta$---the sign is genuinely
distributional, resolved only once the two means are known.
\end{proposition}

\begin{proof}
The partials are read off \eqref{eq:Phi}; $\phio-\qo>0$ because $\phio=F(\KFo)>F(0)=\qo$
($\KFo>0$). Nonzero partials give the nontrivial dependence. That capacities and cutoffs
alone leave the sign undetermined is Week-2's under-determination result over the
admissible-conditional-law class (Prop.~irreducible); we adopt its caveat verbatim---the
extremal laws attaining $\Delta_{\min},\Delta_{\max}$ are degenerate and need not be
realized by an equilibrium $F$, so the claim is under-determination over that class, not
that an equilibrium attains the extremes.\footnote{Verified in \texttt{verify\_week4.py}:
at the representative optimum, sweeping $\mustar$ alone (with $\mumid$ and the cutoffs at
their equilibrium values) makes $\Phi$ cross zero---$\Phi\in[-3.3,+0.4]$---and sweeping
$\mumid$ alone makes it cross---$\Phi\in[-5.3,+0.5]$; so at this optimum each mean is
individually sign-decisive. Whether the single-mean interval contains zero at
\emph{every} interior optimum we do not prove.}
\end{proof}

The per-mean question---does each mean \emph{alone} flip the sign?---has a clean exact
answer, conditional on a single scalar margin.

\begin{theorem}[Per-mean decisiveness, sharp]\label{thm:permean}
At an interior optimum, where $\Delta=\Phi<0$ (Week-2), $\mustar$ is individually
sign-decisive---sweeping it over $(-\KRs,a)$ with $\mumid$ and the cutoffs held at
equilibrium makes $\operatorname{sign}\Phi$ flip---\emph{if and only if}
\[
  \frac{\CM}{R-c}\,(\mustar+\KRs)\;>\;|\Delta| \tag{$\mathrm{D}_{\ell}$}
\]
(the margin by which raising $\mustar$ to the support edge can lift $\Phi$ exceeds the
deficit $|\Delta|$). Symmetrically, $\mumid$ is individually sign-decisive iff
$(\phio-\qo)(\KFo-\mumid)>|\Delta|$ $(\mathrm{D}_m)$. When $(\mathrm{D}_\ell)$ and
$(\mathrm{D}_m)$ both hold, neither mean is redundant and the coordinate set of $\Phi$ is
minimal.
\end{theorem}
\begin{proof}
$\Phi$ is affine and strictly decreasing in $\mustar$ (slope $-\CM/(R-c)$), so over
$\mustar\in(-\KRs,a)$ it ranges over $(\Phi_{\mathrm{lo}},\Phi_{\mathrm{hi}})$ with
$\Phi_{\mathrm{hi}}=\Delta+\tfrac{\CM}{R-c}(\mustar+\KRs)$ (at the upper edge
$\mustar=-\KRs$) and $\Phi_{\mathrm{lo}}=\Delta-\tfrac{\CM}{R-c}(a-\mustar)<\Delta<0$
(at the lower edge). The interval contains $0$ iff $\Phi_{\mathrm{hi}}>0$, i.e.\ iff
$(\mathrm{D}_\ell)$. The $\mumid$ statement is identical with slope $\phio-\qo>0$ and edge
$\mumid=\KFo$. Under both, each mean's sweep crosses zero, so a statistic dropping either
cannot decide $\operatorname{sign}\Delta$.\footnote{Verified in \texttt{verify\_week4.py}:
at the representative optimum $(\mathrm{D}_\ell)$ and $(\mathrm{D}_m)$ hold with margins
$\Phi_{\mathrm{hi}}=+0.40$ and $+0.49$, so both means are individually decisive there.}
\end{proof}

\begin{remark}[What is now pinned, and what remains]
The per-mean flip is now a \emph{theorem}, exact in both directions, gated on the explicit
scalar margins $(\mathrm{D}_\ell),(\mathrm{D}_m)$---no longer a bare numerical
observation. What is not proved is that $(\mathrm{D}_\ell),(\mathrm{D}_m)$ hold at
\emph{every} interior optimum (they hold wherever the sweep margin beats the carbon
deficit, which is checkable per instance and holds throughout the scan). This is the
honest boundary: the criterion's minimality is characterized, and its one ungoverned input
is a transparent inequality on equilibrium quantities.
\end{remark}

%=====================================================================
\section{The carbon phase diagram}\label{sec:phase}
%=====================================================================

Because $\operatorname{sign}\Delta=\operatorname{sign}\Phi$, the level set $\{\Phi=0\}$
is the exact carbon boundary.

\begin{theorem}[The carbon boundary is an interval, under monotonicity]\label{prop:phase}
The primitive space partitions into $\{\Phi>0\}$ (carbon down), $\{\Phi=0\}$ (neutral),
$\{\Phi<0\}$ (up), with boundary the zero set of \eqref{eq:Phi}. If $\Phi(\CR)$ is
monotone along the equilibrium locus (Condition~\textup{S}: $\tfrac{d}{d\CR}\Phi\le0$),
then the carbon-down set is an interval $\{\CR<\CR^{\dagger}\}$ with a single boundary
$\CR^{\dagger}$---the connectedness of Week~2, read off $\Phi$. Condition~\textup{S} holds
throughout the scanned band ($\Phi:{+}1.9\to{-}2.5$, single sign change).
\end{theorem}

\begin{proof}
The partition and boundary are Theorem~\ref{thm:classify}. A monotone $\Phi(\CR)$ with a
sign change has exactly one zero $\CR^{\dagger}$, and $\{\Phi>0\}=\{\CR<\CR^{\dagger}\}$ is
an interval---this is the content given Condition~S. We do \emph{not} derive Condition~S
in closed form: $\Phi$'s dependence on $\CR$ runs through \emph{both} the capacities
($\KRs,\KFs,\KRo,\KFo$ all move with $\CR$) and both conditional means, so no one-term
derivative signs it, and connectedness for general $F$ is exactly Week-2's Condition~MRSV,
which the team could not derive from primitives---so S is the honest hypothesis, checkable
and verified here.\footnote{Verified in \texttt{verify\_week4.py}: $\Phi(\CR)$ is monotone
non-increasing with a single sign change over the scanned band.}
\end{proof}

\begin{remark}[The naive carbon-down region is the corner]
For the \emph{naive} contract that $\Phi$ classifies, Week-2 (Prop.~regimecarbon) found
$\Delta<0$ at every scanned interior optimum: the region $\{\Phi>0\}$ is essentially the
cheap-renewable \emph{corner} $\KRs\to\infty$ (where $\E[\fossil^{*}]\to0$), not a set of
interior optima. So $\Phi$ classifies a sign that is generically negative under the naive
rule; the boundary $\{\Phi=0\}$ is where the corner takes over, consistent with the
viability threshold $\CR^{\dagger}$ of Week~1, \emph{not} a smooth interior level set. The
positive-carbon region opens up only once the contract is optimized (\S\ref{sec:leverage}
and Week~3).
\end{remark}

%=====================================================================
\section{Closing the frontier leverage condition}\label{sec:leverage}
%=====================================================================

Week~3 reduced global carbon-optimality of the surplus contract over the budget family
$d_\tau(x)=\min\{(-x)_++\tau,\KM,(\KF-x)_+\}$ to the leverage condition $\Lambda(\tau)\le1$,
$\Lambda=[G(K(0))-G(K(\tau))]/\Phi_{\mathrm{ft}}(\tau)$, with $G$ consumer fossil and
$\Phi_{\mathrm{ft}}$ the load's footprint. We advance it toward a theorem.

\begin{theorem}[Leverage decomposition]\label{thm:leverage}
The reduction is exact: $C(\tau)-C(0)=\Phi_{\mathrm{ft}}(\tau)\,(1-\Lambda(\tau))$, so
$\tau=0$ minimizes carbon over the family iff $\Lambda(\tau)\le1$. Differentiating
$C(\tau)=G(\KR(\tau),\KF(\tau))+\Phi_{\mathrm{ft}}(\tau)$ (with $G$ independent of $\KM$),
\[
  \frac{dC}{d\tau}=\frac{d\Phi_{\mathrm{ft}}}{d\tau}
  +\big(1-F(\KF)\big)\,\KF'(\tau)-\Big(\int_0^{\KF}\FKR\Big)\,\KR'(\tau),
\]
the three forces being the load's own footprint, the extra consumer fossil from a larger
$\KF$, and the renewable channel. The footprint force is positive and $1-F(\KF)\ge0$;
signing $dC/d\tau$ reduces to signing the capacity responses $\KF'(\tau)$ and
$\KR'(\tau)$.
\end{theorem}
\begin{proof}
The identity is Week~3; the derivative is Leibniz, with $G_{\KF}=1-F(\KF)$,
$G_{\KR}=-\int_0^{\KF}\FKR$, and $\partial G/\partial\KM=0$. ($\tfrac{d\Phi_{\mathrm{ft}}}{d\tau}$
is the total derivative, including $\Phi_{\mathrm{ft}}$'s dependence on the moving
capacities.)
\end{proof}

\begin{theorem}[Leverage bound under capacity monotonicity]\label{thm:leveragebound}
Suppose along the budget path the fossil budget does not raise renewable capacity and does
not lower fossil capacity: $\KR'(\tau)\le0\le\KF'(\tau)$ (Condition~\textup{M}). Then
$G(K(\tau))\ge G(K(0))$, hence
\[
  C(\tau)-C(0)=\Phi_{\mathrm{ft}}(\tau)+\big[G(K(\tau))-G(K(0))\big]\ge\Phi_{\mathrm{ft}}(\tau)\ge0,
  \qquad\text{so}\qquad \Lambda(\tau)\le0<1,
\]
and $\tau=0$ (surplus-only) is the carbon-optimal member of the family. Condition~\textup{M}
holds across every shape and regime scanned---the budget raises $\KF$ (it must, to serve
the fossil it now draws) and does not raise $\KR$ (fossil dispatch earns the $R-c$ margin
on \emph{fossil}, which renewables do not supply)---so surplus-only is carbon-optimal.
\end{theorem}
\begin{proof}
$G(\KR,\KF)=\int_0^{\KF}(1-F(x;\KR))\,dx$ has $\partial G/\partial\KF=1-F(\KF)\ge0$ and
$\partial G/\partial\KR=-\int_0^{\KF}\FKR<0$. Under Condition~\textup{M},
$\KF(\tau)\ge\KF(0)$ and $\KR(\tau)\le\KR(0)$, so $G$ moves up on both counts:
$G(K(\tau))\ge G(K(0))$. Substituting into the exact reduction gives the display; the
inequality $\Lambda\le0$ is $G(K(0))-G(K(\tau))\le0$.\footnote{Verified in
\texttt{verify\_week4.py}: along $\tau$, $\KF$ is non-decreasing ($3.4\to5.3$) and $\KR$
is non-increasing, so Condition~\textup{M} holds; $\Lambda(\tau)<0$ throughout.}
\end{proof}

\begin{remark}[What is now proved, and what remains]
The leverage question is closed \emph{conditionally}: $\Lambda\le1$ is a
\emph{theorem} given the capacity-monotonicity Condition~M, and M is a transparent,
checkable statement about the two capacity responses that holds throughout the scanned
region. What is not proved is Condition~M itself in closed form---signing
$\KR'(\tau),\KF'(\tau)$ in the coupled, submodular three-capacity program requires the full
inverse Hessian (Week-2's cofactor identity), which we do not carry out. So the residual
open step is a clean, isolated comparative-statics sign, not the global optimality claim.
\end{remark}

%---------------------------------------------------------------------
\section{Numerical validation}\label{sec:num}

\texttt{verify\_week4.py} checks, against adaptive quadrature on Beta scale mixtures
(symmetric, left- and right-skewed, mildly skewed): the classification identity
$\Phi=\Delta$ to $\sim10^{-9}$ and $\operatorname{sign}\Delta=\operatorname{sign}\Phi$ in
every cell (Theorem~\ref{thm:classify}); the partial signs
$\partial\Phi/\partial\mustar<0$, $\partial\Phi/\partial\mumid>0$
(Proposition~\ref{thm:minimal}); the per-mean margins $\mathrm{D}_\ell,\mathrm{D}_m$
($\Phi_{\mathrm{hi}}=+0.40,+0.49>0$) and that the exact \emph{iff}
$\mathrm{D}_\ell\Leftrightarrow$ (sweep flips) holds (Theorem~\ref{thm:permean});
monotonicity of $\Phi(\CR)$, Condition~S (Theorem~\ref{prop:phase}); and Condition~M
($\KF$ non-decreasing, $\KR$ non-increasing along $\tau$) with the resulting
$\Lambda(\tau)<0$ (Theorem~\ref{thm:leveragebound}). The suite reports a single pass count.

\end{document}
